\chapter{Experimentation}\label{chap:experimentation}

% --------------------- Section 1 ---------------------

\section{Introduction}\label{sec:experimentation-introduction}
To validate the utility of stable specialization, we conducted a series of experiments focusing on ecosystem analysis.
These experiments aimed to identify real-world scenarios where stable specialization could provide significant benefits, particularly in terms of performance and code maintainability.
The experiments were designed to analyze a diverse set of public Rust codebases, examining their use of traits and generic programming patterns.
By applying our specialization technique to these codebases, we sought to demonstrate its practical applicability and effectiveness in real-world scenarios.

\paragraph{Project Selection}
The first step to conduct this research was to select a set of representative Rust projects from the open-source ecosystem.
For this purpose, we selected some of the most downloaded crates from \texttt{crates.io} as well as some popular projects and other less known ones.

\paragraph{Pattern Identification}
To keep the analysis focused and feasible, we concentrated on identifying specific patterns that are likely to benefit from stable specialization.
These patterns include:
\begin{itemize}
	\item \textbf{Manual Dispatch}: In the trait we have a version of the function for each type we want to support and in the implementation we manually dispatch to the correct version based on the type of the input.
	\item \textbf{Monomorphized Traits}: A trait is defined for each type we want to support and each implementation provides a monomorphized version of the function.
	\item \textbf{Standalone Functions}: A standalone function is defined for each type we want to support and the caller is responsible for calling the correct version based on the type of the input.
	\item \textbf{Inherent Impls}: A struct has inherent implementations for each type we want to support, providing specialized versions of the function.
\end{itemize}
Other than these, we've also pointed out some additional patterns that we did not include in the analysis because they were less common or harder to analyze in a systematic way.
These other patterns include:
\begin{itemize}
	\item \textbf{Type Level Dispatch}: A type-level mechanism is used to select the appropriate implementation at compile time, this does not support default implementations.
	\item \textbf{Macros}: Macros are used to generate specialized implementations based on the input type, this reduces boilerplate but does not support default implementations.
	\item \textbf{Pattern matching blocks}: Pattern matching is used to select the appropriate implementation based on the type of the input.
	\item \textbf{If blocks}: Conditional statements are used to select the appropriate implementation based on the type of the input.
\end{itemize}
Many of these patterns can be seen in the code examples in Listings~\ref{fig:standalone-functions} to~\ref{fig:type-level-dispatch} and the baseline implementations using unstable specialization and \mytool{} specialization can be seen in Listings~\ref{fig:unstable-specialization} and ~\ref{fig:tool-specialization} respectively.
Each pattern represents a different approach to achieving specialization in Rust and our analysis aimed to determine how stable specialization could enhance these patterns.
It's important to note that even though each of these patterns could be implemented using stable specialization, some patterns are more likely to benefit from it than others.
Anyways, we believe that stable specialization could provide significant benefits across all these patterns by reducing boilerplate code, improving performance and enhancing code readability and maintainability.
Furthermore, many of these patterns do not support certain specialization features, like specialization based purely on trait bounds or default implementations.
By leveraging stable specialization, developers can achieve more efficient and maintainable code across a variety of programming patterns.

\begin{figure}[htbp!]
	\centering
	\begin{minipage}{0.49\textwidth}
		\centering
		\showrust[0.98\linewidth]{experimentation-unstable.rs}
		\captionsetup{type=listing,skip=2pt}
		\caption{Unstable Specialization}
		\label{fig:unstable-specialization}
	\end{minipage}
	\hfill
	\begin{minipage}{0.49\textwidth}
		\centering
		\showrust[0.98\linewidth]{experimentation-tool.rs}
		\captionsetup{type=listing,skip=2pt}
		\caption{\mytool{} Specialization}
		\label{fig:tool-specialization}
	\end{minipage}
\end{figure}

\begin{figure}[p!]
	\centering
	\begin{minipage}{0.49\textwidth}
		\centering
		\begin{minipage}{\linewidth}
			\centering
			\showrust[0.98\linewidth]{experimentation-standalone-functions.rs}
			\captionsetup{type=listing,skip=2pt}
			\caption{Standalone Functions}
			\label{fig:standalone-functions}
		\end{minipage}
		\vspace{1em}
		\begin{minipage}{\linewidth}
			\centering
			\showrust[0.98\linewidth]{experimentation-inherent-impls.rs}
			\captionsetup{type=listing,skip=2pt}
			\caption{Inherent Impls}
			\label{fig:inherent-impls}
		\end{minipage}
		\vspace{1em}
		\begin{minipage}{\linewidth}
			\centering
			\showrust[0.98\linewidth]{experimentation-manual-dispatch.rs}
			\captionsetup{type=listing,skip=2pt}
			\caption{Manual Dispatch}
			\label{fig:manual-dispatch}
		\end{minipage}
	\end{minipage}
	\hfill
	\begin{minipage}{0.49\textwidth}
		\centering
		\begin{minipage}{\linewidth}
			\centering
			\showrust[0.98\linewidth]{experimentation-monomorphized-traits.rs}
			\captionsetup{type=listing,skip=2pt}
			\caption{Monomorphized Traits}
			\label{fig:monomorphized-traits}
		\end{minipage}
		\vspace{1em}
		\begin{minipage}{\linewidth}
			\centering
			\showrust[0.98\linewidth]{experimentation-type-level-dispatch.rs}
			\captionsetup{type=listing,skip=2pt}
			\caption{Type Level Dispatch}
			\label{fig:type-level-dispatch}
		\end{minipage}
	\end{minipage}
\end{figure}

\paragraph{Limitations}
It's important to note that our analysis was based on heuristics to identify the relevant patterns within the selected codebases.
As such, there is a possibility that some patterns were overlooked or misclassified, which could affect the comprehensiveness of our findings.

% As expected the major results are in projects that heavily utilize traits and generics, as these are the areas where stable specialization is expected to have the most impact.